\section{Experimental Evaluation}

We designed a comprehensive evaluation methodology to assess both the technical performance of ExplainableCropAI and its practical utility for agricultural stakeholders. Our evaluation approach encompasses quantitative metrics for each system component, as well as a controlled user study to measure impact on trust and decision quality.

\subsection{Evaluation Metrics}

\subsubsection{Disease Detection Performance}

To establish a baseline for comparison, we evaluated the disease detection accuracy of our model against several state-of-the-art approaches. For each model, we measured:

\begin{itemize}
    \item \textbf{Accuracy}: Overall percentage of correctly classified images
    \item \textbf{Precision, Recall, and F1-score}: Class-specific and macro-averaged metrics
    \item \textbf{Confusion Matrix}: To identify common misclassification patterns
    \item \textbf{Inference Time}: Average processing time per image on standard hardware
\end{itemize}

We reported performance on both the standard test set and the challenging test set to assess robustness under difficult conditions.

\subsubsection{Explanation Quality}

For visual explanations, we evaluated both objective faithfulness and subjective interpretability:

\begin{itemize}
    \item \textbf{Localization Accuracy}: Intersection over Union (IoU) between explanation heatmaps and ground truth disease regions
    \item \textbf{Insertion and Deletion Scores} \cite{petsiuk2018rise}: Measuring how explanation-guided image perturbation affects model output
    \item \textbf{Sparsity}: Percentage of the image highlighted by the explanation (lower is better, indicating more focused explanations)
\end{itemize}

For natural language explanations, we assessed:

\begin{itemize}
    \item \textbf{Content Coverage}: Percentage of expected disease-specific information included in the explanation
    \item \textbf{BLEU Score}: Comparing generated explanations with expert-written references
    \item \textbf{Reading Level}: Flesch-Kincaid readability score to ensure accessibility
\end{itemize}

\subsubsection{Uncertainty Estimation}

To evaluate uncertainty quantification, we measured:

\begin{itemize}
    \item \textbf{Calibration Error}: Difference between predicted probability and empirical accuracy
    \item \textbf{Expected Calibration Error (ECE)}: Weighted average of calibration errors across probability bins
    \item \textbf{Uncertainty-Accuracy Correlation}: How well uncertainty estimates correlate with prediction errors
    \item \textbf{Out-of-Distribution Detection}: Ability to assign high uncertainty to images from unseen disease classes
\end{itemize}

\subsubsection{Recommendation Quality}

We evaluated recommendations through expert assessment:

\begin{itemize}
    \item \textbf{Appropriateness}: Rating by agricultural experts of recommendation suitability (1-5 scale)
    \item \textbf{Specificity}: Degree of context-awareness and customization in recommendations
    \item \textbf{Completeness}: Coverage of necessary treatment, prevention, and monitoring advice
\end{itemize}

\subsection{Comparative Analysis}

To contextualize our results, we compared ExplainableCropAI against three alternative approaches:

\begin{enumerate}
    \item \textbf{Black-box CNN}: A standard ResNet50 model without explainability components
    \item \textbf{LIME-explained CNN}: ResNet50 with post-hoc LIME explanations
    \item \textbf{CAM-only}: ResNet50 with Class Activation Mapping but no textual explanations or recommendations
\end{enumerate}

This comparison allowed us to isolate the impact of different explainability components and assess their individual contributions to overall system performance.

\subsection{User Study Design}

To evaluate the practical impact of ExplainableCropAI, we conducted a controlled user study with 45 participants divided into three groups:

\begin{itemize}
    \item 15 agricultural extension officers with formal training in plant pathology
    \item 15 experienced farmers with at least 5 years of paddy cultivation experience
    \item 15 novice farmers with less than 3 years of paddy cultivation experience
\end{itemize}

\subsubsection{Study Protocol}

The study followed a within-subjects design where each participant interacted with three versions of the system:

\begin{enumerate}
    \item \textbf{Prediction-only}: Showing only the disease classification result and confidence score
    \item \textbf{Visual-explanation}: Showing the prediction with Grad-CAM visualization
    \item \textbf{Full-explanation}: Showing the complete ExplainableCropAI output (visual and textual explanations, recommendations)
\end{enumerate}

For each system version, participants completed the following tasks:

\begin{enumerate}
    \item \textbf{Disease Identification}: Reviewing 10 images and either accepting or correcting the system's diagnosis
    \item \textbf{Treatment Decision}: Selecting appropriate treatment actions for each case
    \item \textbf{Trust Assessment}: Rating trust in the system on a 1-7 Likert scale
    \item \textbf{Workload Assessment}: Completing the NASA Task Load Index (TLX)
\end{enumerate}

To control for learning effects, we counterbalanced the order of system versions and used different image sets of comparable difficulty for each version.

\subsubsection{Quantitative Measures}

We measured the following outcomes:

\begin{itemize}
    \item \textbf{Decision Accuracy}: Percentage of correct treatment decisions compared to expert consensus
    \item \textbf{Decision Time}: Time taken to make decisions
    \item \textbf{Trust}: Self-reported trust scores for each system version
    \item \textbf{Mental Workload}: NASA TLX scores focusing on mental demand and frustration
    \item \textbf{System Usability}: System Usability Scale (SUS) scores
\end{itemize}

\subsubsection{Qualitative Feedback}

In addition to quantitative measures, we collected qualitative feedback through:

\begin{itemize}
    \item \textbf{Semi-structured interviews}: Exploring participants' experiences with each system version
    \item \textbf{Think-aloud protocols}: During a subset of tasks to understand decision processes
    \item \textbf{Open-ended questions}: About explanation preferences and improvement suggestions
\end{itemize}

\subsection{Field Deployment}

To assess real-world applicability, we conducted a limited field deployment with 5 agricultural extension officers who used the system during regular field visits over a two-week period. The officers used a mobile app version of ExplainableCropAI to assist farmers with disease diagnosis and treatment planning.

Data collected during this deployment included:

\begin{itemize}
    \item \textbf{Usage patterns}: Frequency and context of system use
    \item \textbf{Adoption barriers}: Technical and practical challenges encountered
    \item \textbf{Impact observations}: Qualitative assessment of influence on field practices
\end{itemize}

\subsection{Ethical Considerations}

Our evaluation was conducted with approval from the institutional ethics committee. All participants provided informed consent and were compensated for their time. To ensure that participants were not over-reliant on the system, we included explicit disclaimers about system limitations and emphasized that it should supplement rather than replace expert judgment.

We addressed potential biases in the evaluation by:

\begin{itemize}
    \item Ensuring diverse representation in our participant pool (gender, age, education level)
    \item Including disease cases with varying difficulty levels
    \item Using pre-registered analysis plans to avoid selective reporting
    \item Having evaluations conducted by researchers not involved in system development
\end{itemize}

The field deployment included careful monitoring to ensure that any system recommendations did not lead to harmful agricultural practices. Participants were encouraged to seek confirmation from extension services before implementing unfamiliar treatments.
